% BEGIN CR28_RADIAL_CONSTITUTIVE_EVALUATION
\section{Constitutive continuity and evaluation of the gravitational section}
\label{cr28:sec:radial-evaluation}

The longitudinal construction determines \(L\) before using gravitational identification. The composition of rules R1--R8 preserves the velocity section of \eqref{cr28:eq:same-velocity} and the previously generated return action. With
\[
 L=\frac{5759}{23040}\alpha^{16}L_*,
 \qquad \hbar_{\rm ret}=\eta_{\rm ret}S_*,
 \qquad c_{\rm int}=c_*\widehat c,
\]
the result of \eqref{g26:eq:g-rigido-es} is expressed as
\begin{equation}
 G_{\rm ret}
 =\frac{c_*^3L_*^2}{S_*}
   \left(\frac{5759}{23040}\right)^2
   \frac{\alpha^{32}\widehat c^{\,3}}{\eta_{\rm ret}}
 =\frac{L^2}
 {\hbar_{\rm ret}(\mu_{\rm vac}\varepsilon_{\rm vac})^{3/2}}.
 \label{cr28:eq:g-constitutive-evaluation}
\end{equation}
The channels are the same as those of \eqref{eq:vacuum-angular-channels}:
\[
 q_\pm=\exp\!\left[-\frac{\pi}{180}(A_{\rm deg}\pm C^*_{\rm deg})\right],
 \qquad r_\pm=\frac{1-q_\pm^{90}}{1-q_\pm^{120}},
 \qquad \widehat c=(r_+r_-)^{-1}.
\]
The constitutive bases satisfy \(\varepsilon=\varepsilon_*r_+^2\), \(\mu=\mu_*r_-^2\) and \(c_*=(\mu_*\varepsilon_*)^{-1/2}\). Their product proves \(c_{\rm int}=(\mu\varepsilon)^{-1/2}\). If the dimensional coordinate of \(c_{\rm int}\) is declared directly, its basis is recovered as \(c_*=c_{\rm int}/\widehat c\). The constitutive coefficient is not multiplied again: both expressions represent the same section.

In the representation \(L_*=1\,\mathrm m\), \(S_*=10^{-34}\,\mathrm{J\,s}\) and \(c_{\rm int}=299792458\,\mathrm{m\,s^{-1}}\), evaluation of the readers and radial composition gives
\begin{equation}
 \boxed{
 G_{\rm ret}
 =6.674299659414022706367776189\ldots\times10^{-11}
 \,\mathrm{m^3\,kg^{-1}\,s^{-2}}.}
 \label{cr28:eq:g-evaluation}
\end{equation}
The comparison value \(G_{\rm CODATA\,2022}=6.67430(15)\times10^{-11}\,
\mathrm{m^3\,kg^{-1}\,s^{-2}}\) retains its experimental uncertainty. In the same units, the difference between the evaluation and the central value is \(-3.40585977\ldots\times10^{-18}\), approximately \(-0.00227057\) standard uncertainties. The evaluation digits correspond to the declared readers and bases; uncertainty belongs to the measured result. No reference digit selects \(N\), \(r_N\), the longitudinal rule or the radial norm.

\subsection{Chronological realization of the same constructed radius}
\label{cr28:sec:radial-clock} With the elementary duration
\begin{equation}
 t_0=\frac{\pi L}{54c_{\rm int}},
 \qquad \ell_0=c_{\rm int}t_0=\frac{\pi L}{54},
 \label{cr28:eq:clock-from-length}
\end{equation}
the calendar of \(108\) positions satisfies
\[
 T_{108}=108t_0=\frac{2\pi L}{c_{\rm int}},\qquad
 \omega_{108}=\frac{c_{\rm int}}{L},\qquad
 R_{108}=\frac{c_{\rm int}}{\omega_{108}}=L.
\]
By substitution,
\begin{equation}
 \left(\frac{54}{\pi}\right)^2
 \frac{c_{\rm int}^5t_0^2}{\hbar_a}
 =\frac{c_{\rm int}^3L^2}{\hbar_a}.
 \label{cr28:eq:circular-radial-agreement}
\end{equation}
This identity places the circular selector retained below in the same longitudinal normalization. Its uniqueness theorem preserves its stated geometric premise; the preceding construction additionally specifies the length used in that realization.

The positive character \(Y=I\) gives simultaneously \(R_X=\Lambda_X=L\) and \(M_X=\hbar_a/(c_{\rm int}L)\). For a general character \(Y>0\), the lengths are reciprocal: \(R_X=LY\), \(\Lambda_X=LY^{-1}\). Their product remains \(L^2I\) without their individual values necessarily coinciding. Thus identification of the four lengths in the circular case corresponds to a precise sector of the radial proof and does not eliminate dependence on the character.

The proportionality \(R_X=(G/c_{\rm int}^{2})M_X\) is common to all characters in the domain. Its preservation under change of basis, refinement, minimization with memory, Schur complementation and dynamic normalization has been proved in \cref{g26:sec:rigidez-radial-es}. The five-component carrier and its reduction to two helicities, developed subsequently, describe the tensor representation; they do not replace the proof of the common radial coefficient.
% END CR28_RADIAL_CONSTITUTIVE_EVALUATION
